\section{The Zero-Mesh Reading}

The split leaves a question of scale. If geometry generates gauge and the
finite carrier keeps the two faces apart, how can the electromagnetic
reading be related to the geometric reading without forming the forbidden
interior sum? The answer is Cauchy convergence. The electromagnetic reading
is the zero-mesh limit of the geometric reading: not a completed continuum
assumed in advance, but the limiting gauge face forced as geometric
residuals are refined toward zero.

Begin with the geometric ledger. It carries a mesh, boundary separations,
and residuals that measure how much geometric refinement has not yet been
absorbed. A coarse mesh reads a coarse boundary grammar. A finer mesh can
split cells, add anchors, and test narrower supports. At each stage the
geometric reading is finite. It has not inspected every point. It has only
committed the distinctions its mesh can carry.

The gauge reading is what remains stable under the refinement of that mesh.
As the geometric residual is bounded by a visible \(\epsilon\), and as
\(\epsilon\) is forced downward, the transported orientation becomes less
dependent on the coarse presentation. The charge reading does not wait for a
completed continuum to arrive from outside. It is the value that every
admissible refinement is being squeezed toward. In this sense the
electromagnetic reading is the Cauchy limit of the geometric reading.

The shape is the same as the earlier finite approach to real value:
\[
  0 \leq R_n \leq \epsilon_n,
  \qquad \epsilon_n \to 0 .
\]
Here \(R_n\) is the unread geometric residue of the \(n\)-th mesh, and the
gauge face is read when the residual sequence becomes Cauchy under the
allowed refinements. The lower bound forbids hidden negative debt from
cancelling a positive residue. The upper bound records the mesh error still
being carried. The vanishing of the upper bound forces the limiting reading.

This is what the chapter means by saying that QED is the Cauchy limit of GR.
The phrase should not be inflated into a claim that every smooth physical
theory has been derived from every other. It names a specific grammatical
relation. The geometric reading carries mesh, boundary, and curvature-like
residue. The electromagnetic reading carries the gauge sign that remains as
the geometric mesh is forced to zero. The two are related by a Cauchy
process, not by an interior direct sum.

The Precision example gives the right caution. A spline completion may
assign determinate values between anchors, but those intermediate values are
not new measurements. They are consequences of the admissible completion.
Likewise, the zero-mesh gauge reading is not an observed point inside a
completed continuum. It is the compelled reading of a refinement process
whose residuals have been bounded and forced downward. Precision is a
consequence of coherence; accuracy remains answerable to future records.

This section also explains why the split was necessary. If the geometric and
gauge sectors could be added in the interior, the limit would be too cheap.
One would simply place the electromagnetic term beside the geometric term
and call the result unified. The grammar instead requires a refinement
history. Geometry carries the mesh. Gauge carries the transported sign. The
limit identifies the stable sign as mesh residue vanishes. No interior mixed
scalar is introduced.

The zero-mesh reading is therefore an exterior discipline on a finite
process. Each mesh stage is finite. Each residual bound is finite. Each
extension must preserve previous commitments or explicitly refine them. The
Cauchy property says that after sufficient refinement, later stages cannot
be distinguished by the permitted residual tests beyond the chosen bound.
The electromagnetic reading is the quotient of those late-stage readings
that no admissible test can tell apart.

This is also why baseline dependence survives the limit. Refinement does
not erase orientation. It squeezes geometric residue while preserving the
baseline relation that makes sign readable. A flat-baseline refinement tower
continues to read the electron sign. A tilted-baseline tower continues to
read the positron sign. The limit does not abolish the grammar of
orientation; it makes the transported orientation independent of mesh
presentation.

The chapter has now moved from charged signs to boundary-generated gauge,
from gauge to split carriers, and from split carriers to a limiting
relation. One obligation remains. The Cauchy limit requires a disciplined
family of finite conditions. The grammar must say how those conditions are
ordered, extended, made compatible, and forced to decide the continuum
reading. That is the role of finite Cohen forcing.

